\chapter{Simulation Methodology and Experimental Setup}
\label{chap:simulation_setup}

\section{Introduction}

This chapter describes the corrected simulation methodology used to evaluate the
proposed temporal-state-aware quantum-routing method. The implementation follows
the system model and routing constraints defined in Chapters~3 and~4. In
particular, the simulator explicitly represents physical links, stored
entanglement links, memory decoherence, predicted consumption-time fidelity,
remaining coherence budget, request-level quality-of-service requirements,
probabilistic entanglement generation, swapping outcomes, and budget-limited
regeneration.

All reported experiments use matched random inputs across the compared methods.
For a given seed, the methods receive the same physical topology, initial link
parameters, traffic requests, entanglement-generation outcomes, and
request-service random numbers. This paired design reduces confounding variation
and permits seed-wise statistical comparison.

\section{Simulation Environment}

The simulator was implemented in Python~3.13.13. Network topology construction
and path computation were implemented using NetworkX. NumPy and pandas were used
for numerical processing and tabular analysis, Matplotlib was used for figure
generation, and SciPy was used for paired statistical tests.

The implementation was divided into the following modules:

\begin{itemize}
    \item \texttt{topology.py}: spatial physical-network generation;
    \item \texttt{physical\_links.py}: physical and quantum-link parameters;
    \item \texttt{overlay.py}: physical and artificial entanglement overlay;
    \item \texttt{temporal\_state.py}: age, decoherence, and coherence-budget updates;
    \item \texttt{routing.py}: composite routing cost and route-feasibility checks;
    \item \texttt{traffic.py}: Poisson traffic and heterogeneous request classes;
    \item \texttt{entanglement\_generation.py}: continuous probabilistic ebit generation;
    \item \texttt{regeneration.py}: budget-limited probabilistic regeneration;
    \item \texttt{request\_service.py}: request outcomes and resource consumption;
    \item \texttt{metrics.py}: end-to-end fidelity and performance metrics;
    \item \texttt{baselines.py}: matched benchmark-routing implementations.
\end{itemize}

The source code and experiment runners were maintained under Git version control.
The corrected implementation was developed on the
\texttt{thesis-aligned-simulation} branch so that the preliminary prototype
remained preserved.

\section{Spatial Physical-Network Model}

\subsection{Node Placement}

For a network containing \(N\) nodes, each node was placed uniformly in a
two-dimensional \(50~\mathrm{km}\times50~\mathrm{km}\) region. If node \(i\) has
coordinate \((x_i,y_i)\), the physical distance between nodes \(i\) and \(j\) is

\begin{equation}
d_{ij}
=
\sqrt{(x_i-x_j)^2+(y_i-y_j)^2}.
\end{equation}

Unlike the preliminary prototype, link distance was not sampled independently
after topology generation. The Euclidean distance determined by the node
coordinates was preserved throughout the simulation.

\subsection{Connectivity and Local Links}

A minimum spanning tree over the complete distance-weighted graph was first
constructed to guarantee physical connectivity. Each node was then connected to
its three nearest neighbors. Duplicate undirected edges were automatically
removed. Consequently, every generated physical network was connected while
remaining sparse.

No artificial quantum shortcut was inserted directly into the physical graph.
Physical connectivity and stored-entanglement connectivity were maintained as
separate layers.

\section{Physical-Link Parameters}

For every physical edge \((i,j)\), the elementary entanglement-generation
probability was modeled as

\begin{equation}
p_{ij}^{\mathrm{gen}}
=
p_0
\exp(-\alpha d_{ij}),
\end{equation}

where \(p_0=0.90\) and
\(\alpha=0.02~\mathrm{km}^{-1}\). The expected generation delay was

\begin{equation}
D_{ij}^{\mathrm{gen}}
=
\frac{\tau_{\mathrm{attempt}}}
{\max(p_{ij}^{\mathrm{gen}},\varepsilon)},
\end{equation}

where \(\tau_{\mathrm{attempt}}=0.01~\mathrm{s}\) and
\(\varepsilon=10^{-9}\).

The initial elementary-link fidelity was sampled uniformly from
\([0.90,0.99]\). Unless coherence time was the controlled independent variable,
the memory coherence time was sampled uniformly from
\([0.5,5.0]~\mathrm{s}\). The default link-memory capacity was ten ebits.

The swapping-success probability was set to \(0.95\), the operational error rate
to \(0.01\), the per-swap delay to \(0.002~\mathrm{s}\), and the classical
coordination delay to \(0.005~\mathrm{s}\).

\section{Entanglement Overlay Construction}

\subsection{Physical Overlay Edges}

Each physical edge was copied into the entanglement overlay with a valid
two-node physical path \([i,j]\). A physical overlay edge initially stored five
ebits. Therefore, with a memory capacity of ten, its initial free-memory value
was five.

\subsection{Artificial Entanglement Shortcuts}

Candidate artificial edges were selected only between node pairs that were not
directly connected in the physical graph. For each candidate, the underlying
minimum-distance physical path was computed. Thus, every artificial edge was
backed by a valid sequence of physical edges.

The number of artificial shortcuts was

\begin{equation}
M_{\mathrm{A}}
=
\min
\left(
|\mathcal{C}|,
\operatorname{round}(0.35N)
\right),
\end{equation}

where \(\mathcal{C}\) is the candidate set. Candidates with larger physical hop
count and distance were considered first.

For a physical path containing \(h\) elementary edges, the artificial-link
generation probability was

\begin{equation}
p_{ij}^{\mathrm{A}}
=
\left(
\prod_{\ell=1}^{h}
p_{\ell}^{\mathrm{gen}}
\right)
p_{\mathrm{swap}}^{h-1}.
\end{equation}

The artificial-link delay included the elementary generation delays, swapping
delays, and classical coordination delay. Its coherence time and memory capacity
were bounded by the minimum values along the supporting physical path.

The end-to-end initial fidelity was obtained through Werner-state visibility
composition. For elementary fidelity \(F_\ell\), the corresponding visibility is

\begin{equation}
v_\ell
=
\max
\left(
0,
\min
\left(
1,
\frac{4F_\ell-1}{3}
\right)
\right).
\end{equation}

The artificial-link fidelity was calculated as

\begin{equation}
F_{ij}^{\mathrm{A},0}
=
\frac{1}{4}
+
\frac{3}{4}
\left(
\prod_{\ell=1}^{h}v_\ell
\right)
(1-e_{\mathrm{op}})^{h-1}.
\end{equation}

Each artificial edge initially stored one ebit, leaving nine free memory units
under the default capacity. Resource conservation was enforced as

\begin{equation}
n_{ij}^{\mathrm{ebit}}
+
m_{ij}^{\mathrm{free}}
=
m_{ij}^{\mathrm{cap}}.
\end{equation}

\section{Temporal-State Evolution}

The slot duration was \(\Delta=0.1~\mathrm{s}\). At each slot, the ebit age was
increased and effective fidelity was updated using the Werner-state memory-decay
model

\begin{equation}
F_{ij}^{\mathrm{eff}}(t)
=
\frac{1}{4}
+
\left(
F_{ij}^{0}-\frac{1}{4}
\right)
\exp
\left(
-\frac{A_{ij}(t)\Delta}{T_{2,ij}}
\right).
\end{equation}

The fidelity predicted at the expected consumption time was

\begin{equation}
\widehat{F}_{ij}^{\mathrm{cons}}(t)
=
\frac{1}{4}
+
\left(
F_{ij}^{\mathrm{eff}}(t)-\frac{1}{4}
\right)
\exp
\left(
-\frac{\widehat{D}_{ij}^{\mathrm{use}}(t)}
{T_{2,ij}}
\right).
\end{equation}

The remaining coherence budget was

\begin{equation}
B_{ij}(t)
=
T_{2,ij}
-
A_{ij}(t)\Delta
-
\widehat{D}_{ij}^{\mathrm{use}}(t).
\end{equation}

The default minimum link fidelity was \(0.70\). A link was unusable when its
predicted consumption fidelity fell below this threshold, its coherence budget
was non-positive, or it contained no stored ebit. Free memory was not treated as
a service-feasibility constraint because zero free memory means that the memory
is full, not that its stored ebits are unavailable.

The normalized safe and warning coherence-budget thresholds were \(0.30\) and
\(0.10\), respectively.

\section{Traffic and Request Model}

\subsection{Poisson Arrivals}

The number of requests arriving in a slot was sampled from a Poisson
distribution. Unless otherwise stated, the mean arrival rate was one request per
slot. Source and destination nodes were sampled uniformly without replacement.

\subsection{Request Classes}

Three service classes were used:

\begin{table}[htbp]
\centering
\caption{Entanglement-request service classes.}
\label{tab:request_classes}
\begin{tabular}{lccc}
\hline
Service class & Probability & Required fidelity & Deadline \\
\hline
Standard & 0.50 & 0.80 & 5 slots \\
High fidelity & 0.30 & 0.85 & 4 slots \\
Urgent & 0.20 & 0.75 & 2 slots \\
\hline
\end{tabular}
\end{table}

Every request demanded one end-to-end entangled pair. The deadline in seconds
was obtained by multiplying the slot deadline by \(\Delta\).

\section{Entanglement Generation, Service, and Regeneration}

\subsection{Continuous Ebit Generation}

At every slot, each overlay edge with free memory performed at most one
generation attempt. The attempt succeeded according to the edge-specific
generation probability. A successful attempt added one stored ebit, reduced free
memory by one, reset the ebit age, and restored the initial fidelity and
coherence budget.

\subsection{Request Service}

A request was accepted only when all of the following conditions held:

\begin{enumerate}
    \item a usable route existed;
    \item every route edge contained the demanded stored ebit;
    \item every route edge met the minimum link-fidelity constraint;
    \item route-level predicted fidelity met the request requirement;
    \item expected completion delay did not exceed the deadline;
    \item the probabilistic swapping operation succeeded.
\end{enumerate}

Following successful service, one stored ebit was consumed from every overlay
edge on the route and the corresponding memory was released. Therefore, finding
a path alone was not counted as request acceptance.

\subsection{Budget-Limited Regeneration}

Warning, critical, and expired links were eligible for regeneration. Candidate
priority combined temporal urgency, fidelity deficit, coherence-budget risk,
structural importance, blocked-request count, success probability, and expected
generation delay. The priority was proportional to expected regeneration
utility rather than urgency alone.

Each regeneration attempt succeeded probabilistically. A successful attempt
reset age, restored fidelity, created a fresh stored ebit, and recorded the
regeneration delay. Attempt, success, and failure counts were collected
separately.

\section{Routing Methods}

Three methods were compared in the main baseline experiment:

\begin{enumerate}
    \item \textbf{Shortest Path}: minimum-hop routing over currently usable
    overlay edges;
    \item \textbf{Fidelity-Based}: routing that minimizes the negative logarithm
    of composed Werner-state visibility;
    \item \textbf{Proposed}: the temporal-state-aware composite-cost method with
    route-level fidelity and delay feasibility checks and proactive
    regeneration.
\end{enumerate}

All methods used identical link-availability conditions and matched stochastic
inputs. Independent graph copies were used so that one method's resource
consumption did not modify another method's state.

The proposed composite cost used the following weights:

\begin{table}[htbp]
\centering
\caption{Composite routing-cost weights.}
\label{tab:routing_weights}
\begin{tabular}{lc}
\hline
Cost component & Weight \\
\hline
Predicted-fidelity risk & 0.30 \\
Age risk & 0.15 \\
Coherence-budget risk & 0.25 \\
Regeneration risk & 0.15 \\
Structural risk & 0.15 \\
\hline
\end{tabular}
\end{table}

At most the 20 lowest-cost simple candidate routes were checked. This bound
prevented uncontrolled path enumeration in denser networks while preserving
alternative-route feasibility evaluation.

\section{Experimental Scenarios}

The corrected evaluation included six experiments:

\begin{enumerate}
    \item \textbf{Baseline comparison}: Shortest Path, Fidelity-Based, and
    Proposed routing over 20 seeds, 20 nodes, and 100 slots.
    \item \textbf{Multi-seed comparison}: paired Proposed-versus-Static
    evaluation over 20 seeds.
    \item \textbf{Ablation study}: Static, Adaptive-Only, and Proposed variants
    over 20 seeds.
    \item \textbf{Coherence-time analysis}: controlled
    \(T_2\in\{0.5,1,2,5,10,20\}~\mathrm{s}\), with 20 seeds per setting.
    \item \textbf{Regeneration-budget analysis}: budgets
    \(\{0,1,3,5,10\}\) attempts per slot, with 20 seeds per setting.
    \item \textbf{Network scalability}: network sizes
    \(\{6,20,40,60,100\}\), with 20 seeds per size.
\end{enumerate}

The ablation, coherence-time, regeneration-budget, multi-seed, and scalability
experiments used 100 slots per run. The default network size for the multi-seed,
ablation, coherence-time, and regeneration-budget experiments was 50 nodes.

\section{Evaluation Metrics}

The principal metrics were:

\begin{itemize}
    \item request acceptance ratio;
    \item accepted and failed request counts;
    \item average accepted-route fidelity;
    \item average accepted-route completion delay;
    \item average accepted-route hop count;
    \item regeneration attempts, successes, and failures;
    \item physical and overlay edge counts;
    \item total runtime and runtime per request.
\end{itemize}

Fidelity, delay, and hop metrics were averaged over accepted requests. When
combining multiple seeds, weighted means used accepted-request counts as weights
to avoid bias from seeds with zero accepted requests.

\section{Statistical Methodology}

Every reported configuration was evaluated using 20 independent seeds. Mean
acceptance ratios were reported with two-sided 95\% confidence intervals.
Because methods received matched inputs for each seed, paired two-sided
\(t\)-tests were used for method comparisons. The significance threshold was
\(p<0.05\), while results with \(p<0.001\) were treated as highly significant.

\section{Implementation Validation}

Eight automated tests verified the core scientific and resource invariants:

\begin{enumerate}
    \item physical topology connectivity;
    \item preservation of spatial link distance;
    \item validity of artificial-link physical paths;
    \item agreement with the Werner-state decay equation;
    \item decreasing fidelity with additional swapping hops;
    \item request-service memory and ebit conservation;
    \item generation-capacity conservation;
    \item reproducibility of seeded traffic.
\end{enumerate}

All eight tests passed before the full experiments were executed.

\section{Chapter Summary}

This chapter presented the corrected and reproducible simulation methodology.
The evaluation separates physical and artificial connectivity, applies
consumption-time temporal feasibility, models stochastic generation and
regeneration, consumes resources after successful service, and compares methods
using matched inputs and paired statistics. The resulting experimental findings
are presented in Chapter~6.